% !PRINTED 4
\weberpage{4}
%<<CONT
\gesperrt{the general elliptic differential.} But this general integral and differential can now be reduced to essentially simpler ones. First, $\Phi$ can be put in the form:
\[
\frac{A + B\sqrt{f(x)}}{C + D\sqrt{f(x)}},
\]
where $A$, $B$, $C$, $D$ are integral rational functions of $x$, and, on multiplying numerator and denominator of this fraction by $C - D\sqrt{f(x)}$, in the form
\[
\Psi(x) + \frac{\Phi(x)}{\sqrt{f(x)}},
\]
where $\Psi(x)$ and $\Phi(x)$ are integral or fractional rational functions of $x$, namely
\[
\Psi(x) = \frac{AC - BDf(x)}{C^2 - D^2f(x)},
\]
\[
\Phi(x) = \frac{(BC - AD)f(x)}{C^2 - D^2f(x)}.
\]

We now leave out of account the rational integral
\[
\int \Psi(x)\, dx,
\]
which leads to logarithms and algebraic functions, and occupy ourselves only with the elliptic integral
\weq{2}{\int \frac{\Phi(x)\, dx}{\sqrt{f(x)}}}
or with the corresponding elliptic differential
\weq{3}{\Phi(x) \frac{dx}{\sqrt{f(x)}}.}

It is sometimes useful to consider this differential in the homogeneous form; to this end we put for $x$ the ratio of two variables $x:y$, and accordingly for $dx$
\[
\frac{y\, dx - x\, dy}{y^2}.
\]
If then
\[
f(x, y) = a_0 x^4 + 4 a_1 x^3 y + 6 a_2 x^2 y^2 + 4 a_3 x y^3 + a_4 y^4,
\]
and $\Phi(x, y)$ is a homogeneous function of order 0, our elliptic differential assumes the shape
\weq{4}{\Phi(x, y) \frac{y\, dx - x\, dy}{\sqrt{f(x, y)}}.}
